\chapter*{Conclusion g\'en\'erale et perspectives}
\addcontentsline{toc}{chapter}{Conclusion g\'en\'erale et perspectives}

\guided{La conclusion doit r\'epondre explicitement aux objectifs annonc\'es dans l'introduction. Ne pas introduire de nouveaux r\'esultats. Terminer par des perspectives coh\'erentes, r\'ealistes et argument\'ees.}

Ce projet a permis de traiter la probl\'ematique relative \`a \todo{rappeler le sujet}. Les travaux men\'es ont conduit \`a \todo{r\'esumer les r\'esultats majeurs et les contributions principales}. Les objectifs atteints concernent notamment \todo{lister les acquis principaux}, tandis que certaines limites subsistent, en particulier \todo{indiquer les limites techniques, m\'ethodologiques ou organisationnelles}.

D'un point de vue acad\'emique et professionnel, ce travail met en \'evidence \todo{les comp\'etences mobilis\'ees : analyse, conception, validation, esprit critique, autonomie}. Les perspectives envisag\'ees portent sur \todo{am\'eliorations futures, extension du p\'erim\`etre, industrialisation, passage \`a l'\'echelle, nouveaux jeux de donn\'ees, validation sur site, etc.}.
